\section{Introduction}
\label{sec:introduction}

Long-horizon robotic manipulation executes a sequence of physical interactions whose success depends on the evolving state of the world: a pick-and-place task is complete only if the robot grasps the object, lifts it without dropping it, and places it at the target --- in that order. The harder problem is not an individual step failing but how execution responds afterwards. One failure mode is stagnation: the policy keeps trying to grasp the same orange until the episode times out, even when it is not succeeding and other oranges remain available. A second, more speculative risk is state poisoning: a missed or fumbled grasp can knock an orange into a rarely-seen position, and from such an out-of-distribution state the policy may act unreliably. The latter is treated as a hypothesis rather than an established mechanism.

Learned policies are especially vulnerable to this. Trained on successful demonstrations, their state distribution is dominated by runs in which previous steps went well; during execution, small deviations push them into underrepresented states where they still produce plausible-looking but unreliable actions, and over a long horizon these errors accumulate. Vision-Language-Action (VLA) models add language as an interface for specifying behaviour, but language alone does not solve the problem: an instruction such as ``place the object in the plate'' assumes the object is already held, and if that assumption is false it offers no mechanism to interrupt or correct execution.

This project asks whether recovery can be enabled by changing the execution structure \emph{around} the policy. The hypothesis is that a long-horizon task becomes more robust if it is decomposed into language-conditioned subtasks wrapped by an orchestrator that checks whether each subtask achieved its intended physical outcome and, on failure, recovers at two levels: retrying the same sub-goal (recovering the specific orange) or redirecting to a feasible alternative orange (recovering overall task progress). This is tested in a simulated kitchen task where an SO-101 arm must pick three oranges and place them into a plate.

The contributions are threefold. (i) ACT and SmolVLA are compared in a monolithic full-task setting, showing that changing the policy class alone is not enough to solve the task reliably. (ii) A language-conditioned Teleop subtask dataset with randomised pick orders is built, together with an orchestrator that monitors subtask outcomes, enables local retry, and redirects after a GRASP timeout; because the policy follows targeting prompts only from configurations far from the oranges, each target change is prefaced by a scripted home-pose reset. (iii) The orchestrated system is reused for autonomous data generation, evaluating whether the added successful subtasks improve the policy --- a useful negative finding: under the seeded protocol, the autonomous data reduces orchestrated subtask performance.

\Cref{sec:related} situates this work against prior approaches to VLA manipulation, long-horizon decomposition, and failure detection and recovery; \Cref{sec:methods} describes the simulation environment, dataset construction, policy training, and orchestrator design; \Cref{sec:experiments} reports the experimental results; \Cref{sec:discussion} discusses implications and future work.
